\section{Parser (syntaktický analyzátor)}

Syntaxický analyzátor (parser) je zodpovědný za kontrolu struktury programu a vytvoření jeho mezivýsledného reprezentace \textbf{ve formě abstraktního syntaktického stromu (AST)}. Implementace je založena na hybridním přístupu, který kombinuje dvě metody analýzy:

\begin{itemize}
    \item \textbf{Metoda rekurzivního sestupu (Recursive Descent):} Používá se pro zpracování hlavních konstrukcí jazyka (definice funkcí, cykly, podmíněné příkazy).
    \item \textbf{Precedenční analýza (Precedence Analysis):} Používá se výhradně pro zpracování výrazů (matematických, logických a řetězcových).
\end{itemize}

\subsection{Struktura implementace}

Zdrojový kód modulu je rozdělen do několika souborů podle logických úloh:

\begin{description}
    \item[\texttt{parser.c} (Hlavní modul):]
    Implementuje metodu rekurzivního sestupu na základě LL gramatiky a poskytuje jedinou veřejnou funkci \texttt{parser\_run}, která spouští analýzu a vrací kořen AST programu.

    \item[\texttt{expression.c} (Rozhraní pro výrazy):]
    Slouží jako rozhraní, které propojuje hlavní parser s modulem precedenční analýzy a předává mu zpracování výrazů a tvorbu jejich AST.

    \item[\texttt{precedence.c} (Tabulka priorit):]
    Obsahuje implementaci tabulky priorit a definuje logiku přesunu/redukce pro korektní parsování matematických a logických výrazů.

    \item[\texttt{stack\_precedence.c} (Zásobník analýzy):]
    Implementuje specializovaný zásobník pro algoritmus shift-reduce, který uchovává operátory a uzly AST pro dynamickou konstrukci stromu výrazů.
\end{description}

\subsection{Návrh syntaktického analyzátoru}


\begin{itemize}
    \item \textbf{Volba použití Abstract Syntax Tree (AST):}
    Abychom při generování kódu nemuseli složitě vracet zpět a doplňovat chybějící adresy, používáme AST. Tím se práce rozdělí do dvou kroků: nejprve vytvoříme strom a najdeme všechny funkce, a teprve potom generujeme kód. Tento postup dělá implementaci jednodušší a spolehlivější.

    \item \textbf{Důvod pro hybridní přístup:}
    Pro strukturální část jazyka (podmínky, cykly) jsme zvolili \textbf{rekurzivní sestup}, protože přirozeně kopíruje zanořenou strukturu kódu a je snadno laditelný.
    Naproti tomu pro výrazy je rekurzivní sestup neefektivní kvůli nutnosti řešit priority operátorů (což vede k hlubokému zanoření volání). Proto jsme pro výrazy zvolili \textbf{precedenční analýzu} (bottom-up), která priority zpracovává tabulkově a efektivněji.

\end{itemize}

